% Notes of the Montreal seminar, 1970 (cote 129, typed part): archive cote 129, site page 60 (PDF page 61).
% typescript. Draft compared: none.
% AI-transcribed by Claude Opus 5.5 (claude-opus-5-5, adaptive thinking, effort high), Message Batches API, batch mtl129_full (msgbatch_01AGnDKHwy6ATCM7ehUa6263), 2026-11-13; not yet checked by eye.
\scanpage{129/60}
\mspage{\hand{III-5}}

et donc (fidèlement) plate.

Par le critère de platitude fibre à fibre (EGA, IV, 11.3.1)
on en déduit que $G(i)$ est plat sur $S$. Alors $G(i)$ est plat, fini,
localement de présentation finie sur $S$, et est donc localement libre.

Inversement, si $G(1)$ est représentable par un schéma fini
localement libre sur $S$ et si $G(n)$ est $\Lambda_n$-plat, alors $G(n)$ admet
une suite de décomposition finie $p^iG(n)/p^{i+1}G(n) \approx G(1)$ et $G(n)$ est
représentable par un schéma fini localement libre (II, 1.6). On a donc:

2.4. \uline{Proposition}: Soit $G(n)$ un groupe annulé par $p^n$ et plat sur
$\Lambda_n$. Alors $G(n)$ est représentable par un schéma fini localement libre
sur $S$ si et seulement si l'un des $G(i)$, $1 \leq i \leq n$, l'est.

3. \uline{Groupes de Barsotti-Tate tronqués d'échelon $n$}

3.1. Nous allons voir que la définition d'un groupe de Barsotti-Tate
tronqué d'échelon $n$ fait intervenir la platitude sur $\Lambda_n$. Pour $n = 1$
cette condition est vide et il est nécessaire d'introduire une condition
supplémentaire que l'on justifiera a posteriori (3.3).

Soit $G$ un schéma en groupes (commutatifs) plat sur $S$. On
considère le $\mathbb{F}_p$-schéma fermé $S_o \subset S$ défini par l'idéal $p\underline{\mathcal{O}}_S$ et soit
$G_o$ le schéma sur $S_o$ obtenu par changement de base. On peut alors cons-%
truire les morphismes $\mathfrak{f}_{G_o/S_o}$ et $\mathfrak{v}_{G_o}$. Si $G_o$ est un $\Lambda_1$-module, on a
$\mathfrak{f}_{G_o/S_o} \circ \mathfrak{v}_{G_o} = 0$.

\mspage{\hand{60}}
